\chapter{The Spectral Theorem and Its Relatives}\label{linear:ch:spectral}

\begin{center}\itshape
I developed my theory of infinitely many variables from purely mathematical interests, and even called it ``spectral
analysis'' without any presentiment that it would later find application to the actual spectrum of physics.
\upshape(attributed to David Hilbert, in the 1920s)
\end{center}

In 1904 David Hilbert, studying integral equations, needed a name for the set of values $\lambda$ for which $(A-\lambda I)$
fails to be invertible, and he borrowed one from optics: the \emph{spectrum}. It was a metaphor. A prism splits white light
into colours, and Hilbert's theorem split a symmetric operator into simple pieces, each acting by a single number. Twenty
years later, Heisenberg, Born, Jordan and Schr\"odinger discovered that the energy levels of the hydrogen atom, whose
differences are the frequencies of the lines in its actual, physical spectrum, are the eigenvalues of a symmetric
operator. The metaphor had been literally true all along. Hilbert's reaction, quoted above as it was reported, became part
of the folklore of the subject.

Long before Hilbert, the theorem in finite dimension was known: Cauchy (1829) proved that a real symmetric matrix has real
eigenvalues and an orthonormal basis of eigenvectors. What the twentieth century added is the \emph{variational} point of
view. Lord Rayleigh, in \emph{The Theory of Sound} (1877), estimated the lowest frequency of a vibrating string or plate by
minimising a ratio of energies, and Ernst Fischer (1905) and Richard Courant (1920) characterised \emph{every} eigenvalue
as a min-max of such ratios. Hermann Weyl (1912) used this to prove that the eigenvalues of a drum grow in proportion to $n$
divided by its area, and in passing found inequalities between the eigenvalues of $A$, $B$ and $A+B$ that still bear his
name.

This chapter proves the spectral theorem for normal operators, then studies eigenvalues of Hermitian matrices as
optimisation problems: Rayleigh quotients, the Courant--Fischer theorem, Cauchy's interlacing, Weyl's inequalities, and two
theorems that tell exactly which diagonals a symmetric matrix with given eigenvalues can have (Schur and Horn) and how far
eigenvalues move when a matrix moves (Hoffman and Wielandt).

\begin{think}
Is there a real symmetric $3\times3$ matrix with diagonal entries $3,0,0$ and eigenvalues $2,2,-1$? With diagonal entries
$1,1,1$ and eigenvalues $3,0,0$? The traces agree in both cases, so the trace alone cannot decide. Section~\ref{linear:sec:schurhorn}
gives a complete answer; the first question has an answer that fits on one line once you know the Rayleigh quotient.
\end{think}

\section{The spectral theorem}

Throughout, $V$ is a finite-dimensional inner product space over $\C$ (or $\R$, when said), and $A^*$ is the adjoint of
Chapter~\ref{linear:ch:inner}, the conjugate transpose in an orthonormal basis.

\begin{definition}\label{linear:def:normal}
An operator $A$ is \emph{normal} if $AA^*=A^*A$. Special cases: \emph{Hermitian} (or \emph{self-adjoint}; \emph{symmetric}
over $\R$) if $A^*=A$; \emph{skew-Hermitian} if $A^*=-A$; \emph{unitary} if $A^*A=I$.
\end{definition}

\begin{theorem}[spectral theorem]\label{linear:thm:spectral}
An operator on a finite-dimensional complex inner product space is normal if and only if it has an orthonormal basis of
eigenvectors. In matrix terms, $A\in M_n(\C)$ is normal if and only if $A=UDU^*$ with $U$ unitary and $D$ diagonal.
\end{theorem}
\begin{proof}
If $A=UDU^*$, then $A^*=U\bar DU^*$, and $AA^*=U D\bar DU^*=U\bar DDU^*=A^*A$ because diagonal matrices commute.

Conversely, by Schur's theorem (Theorem~\ref{linear:thm:schur}, orthonormal version) there is a unitary $U$ with $T=U^*AU$ upper
triangular, and $T$ is normal: $TT^*=U^*AA^*U=U^*A^*AU=T^*T$. Compare the $(1,1)$ entries of $T^*T$ and $TT^*$: the first is
the squared norm of the first column of $T$, which is $|t_{11}|^2$, and the second is the squared norm of the first row,
$\sum_j|t_{1j}|^2$. Hence $t_{1j}=0$ for $j>1$. Then $T=\operatorname{diag}(t_{11},T')$ with $T'$ upper triangular and
normal, and induction shows that $T$ is diagonal.
\end{proof}

\begin{corollary}\label{linear:cor:hermitian}
The eigenvalues of a Hermitian operator are real, those of a unitary operator have modulus $1$, and eigenvectors of a
normal operator for distinct eigenvalues are orthogonal. A real symmetric matrix is $A=QDQ^{\mathsf T}$ with $Q$ real
orthogonal and $D$ real diagonal.
\end{corollary}
\begin{proof}
For $A=UDU^*$, $A^*=A$ means $\bar D=D$, and $A^*A=I$ means $\bar DD=I$. Eigenvectors for distinct eigenvalues lie in
different eigenspaces, which are spanned by different vectors of the orthonormal eigenbasis. For a real symmetric $A$, the
complex eigenvalues are real, so $\chi_A$ splits over $\R$, and Theorem~\ref{linear:thm:schur} over $\R$ gives a real orthogonal
$Q$ with $Q^{\mathsf T}AQ$ upper triangular; being symmetric, it is diagonal.
\end{proof}

Grouping the eigenvectors by eigenvalue, the theorem says $A=\sum_\lambda\lambda P_\lambda$, where $P_\lambda$ is the
orthogonal projection onto the eigenspace of $\lambda$; these projections are mutually orthogonal and add up to $I$. This is
the \emph{spectral decomposition}, the form that generalises to infinite dimension (Tome~7, Chapter~4). It makes functions
of a normal operator trivial: $f(A)=\sum_\lambda f(\lambda)P_\lambda$, so $e^A$, $|A|$, and, for a Hermitian operator with
eigenvalues $\ge0$, the square root $\sqrt A$ all make sense (Exercise~\ref{linear:ex:sqrtpsd}).

Hermitian operators are the matrices of Hermitian forms, and a Hermitian $A$ with all eigenvalues $\ge0$ (equivalently
$x^*Ax\ge0$ for all $x$, by expanding $x$ in the eigenbasis) is \emph{positive semidefinite}, written $A\succeq0$. From now
on, the eigenvalues of a Hermitian $A$ are listed in decreasing order, $\lambda_1(A)\ge\lambda_2(A)\ge\dots\ge\lambda_n(A)$.

\section{Rayleigh quotients and the min-max theorem}

The \emph{Rayleigh quotient} of a Hermitian $A$ is $R_A(x)=\frac{\langle Ax,x\rangle}{\langle x,x\rangle}$ for $x\ne0$, a real
number. Take an orthonormal eigenbasis $u_1,\dots,u_n$ with $Au_i=\lambda_iu_i$; for $x=\sum_ic_iu_i$,
\begin{equation}\label{linear:eq:rayleigh}
R_A(x)=\frac{\sum_i\lambda_i|c_i|^2}{\sum_i|c_i|^2},
\end{equation}
a weighted average of the eigenvalues. Hence $\lambda_n\le R_A(x)\le\lambda_1$, with equality on the extreme eigenvectors:
\[
\lambda_1(A)=\max_{x\ne0}R_A(x),\qquad\lambda_n(A)=\min_{x\ne0}R_A(x).
\]
Two consequences come for free. First, $\lambda_1$ is a maximum of the linear functions $A\mapsto\langle Ax,x\rangle$ over unit
vectors $x$, hence a \emph{convex} function of $A$: $\lambda_1(A+B)\le\lambda_1(A)+\lambda_1(B)$. Second, each diagonal entry
$a_{ii}=R_A(e_i)$ lies between $\lambda_n$ and $\lambda_1$. That answers the first question of the opening: a diagonal entry
$3$ is impossible when the largest eigenvalue is $2$.

\begin{figure}[ht]
\centering
\fig{linear-ch05-rayleigh}
\caption{The Rayleigh quotient of a $2\times2$ symmetric matrix with eigenvalues $3$ and $1$, on the unit vectors
$u_\theta=(\cos\theta,\sin\theta)$. It is $3\cos^2(\theta-\frac\pi6)+\sin^2(\theta-\frac\pi6)$, maximal on the top eigenvector
and minimal on the other.}
\label{linear:fig:rayleigh}
\end{figure}

What about the eigenvalues in the middle? Each one is a max of mins.

\begin{theorem}[Courant--Fischer]\label{linear:thm:courantfischer}
For a Hermitian operator $A$ on $V$ of dimension $n$ and $1\le k\le n$,
\[
\lambda_k(A)=\max_{\dim W=k}\ \min_{x\in W\setminus0}R_A(x)=\min_{\dim W=n-k+1}\ \max_{x\in W\setminus0}R_A(x).
\]
\end{theorem}
\begin{proof}
For $W_k=\operatorname{span}(u_1,\dots,u_k)$, formula~\eqref{linear:eq:rayleigh} with $c_i=0$ for $i>k$ gives $\min_{W_k}R_A=
\lambda_k$. Now let $W$ be any subspace of dimension $k$. The subspace $U=\operatorname{span}(u_k,\dots,u_n)$ has dimension
$n-k+1$, and $\dim W+\dim U=n+1>n$, so $W\cap U$ contains some $x\ne0$. For that $x$, only $c_k,\dots,c_n$ are non-zero
and $R_A(x)\le\lambda_k$. So $\min_WR_A\le\lambda_k$, which proves the first formula. The second is the first applied to
$-A$, whose $k$-th largest eigenvalue is $-\lambda_{n-k+1}(A)$.
\end{proof}

The point of the theorem is that it characterises $\lambda_k$ without mentioning eigenvectors, so it can be used to compare
the eigenvalues of different operators. The first comparison is between a matrix and its corners.

\begin{theorem}[Cauchy interlacing]\label{linear:thm:interlacing}
Let $A$ be Hermitian of size $n$, and $B=P^*AP$ its compression to a subspace of dimension $m$, where $P\in M_{n,m}$ has
orthonormal columns; for instance, $B$ is a principal $m\times m$ submatrix of $A$. Then for $1\le k\le m$,
\[
\lambda_k(A)\ge\lambda_k(B)\ge\lambda_{k+n-m}(A).
\]
In particular, the eigenvalues $\mu_1\ge\dots\ge\mu_{n-1}$ of an $(n-1)\times(n-1)$ principal submatrix interlace those of
$A$: $\lambda_1\ge\mu_1\ge\lambda_2\ge\mu_2\ge\dots\ge\mu_{n-1}\ge\lambda_n$.
\end{theorem}
\begin{proof}
For $y\in\C^m$, $R_B(y)=\frac{\langle APy,Py\rangle}{\langle Py,Py\rangle}=R_A(Py)$, and $y\mapsto Py$ maps the subspaces of
$\C^m$ bijectively onto the subspaces of $\operatorname{Im}P$, preserving dimension. So by Theorem~\ref{linear:thm:courantfischer},
$\lambda_k(B)$ is the max-min of $R_A$ over $k$-dimensional subspaces of $\operatorname{Im}P$, which is at most the
max-min over all $k$-dimensional subspaces, $\lambda_k(A)$. Similarly $\lambda_k(B)$ is the min-max of $R_A$ over subspaces
of $\operatorname{Im}P$ of dimension $m-k+1$, which is at least the min-max over all subspaces of that dimension, namely
$\lambda_{n-(m-k+1)+1}(A)=\lambda_{k+n-m}(A)$.
\end{proof}

\begin{figure}[ht]
\centering
\fig{linear-ch05-interlace}
\caption{The eigenvalues of the top-left $k\times k$ corners of a random symmetric $6\times6$ matrix, for $k=1,\dots,6$.
Each row interlaces the next one: the thin lines show that every eigenvalue of a corner lies between two consecutive
eigenvalues of the next corner. Such triangular arrays are called Gelfand--Tsetlin patterns.}
\label{linear:fig:interlace}
\end{figure}

This is the theorem promised in Chapter~\ref{linear:ch:inner}: the Jacobi matrix $J_n$ of orthogonal polynomials is the corner of
$J_{n+1}$, so the zeros of consecutive orthogonal polynomials interlace (Figure~\ref{linear:fig:legendre}). In Appendix~H the same
theorem, applied to a cleverly signed adjacency matrix of the cube $\{0,1\}^n$, proves the sensitivity conjecture. Its converse
also holds: any two interlacing sequences are the spectra of a matrix and of a corner of it (Exercise~\ref{linear:ex:bordered}).

\section{Weyl's inequalities}

What can be said about the eigenvalues of $A+B$ when we know those of $A$ and $B$? Not everything, since the answer depends
on how the eigenvectors sit, but a lot.

\begin{theorem}[Weyl, 1912]\label{linear:thm:weyl}
For Hermitian $A,B$ of size $n$ and indices with $i+j-1\le n$,
\[
\lambda_{i+j-1}(A+B)\le\lambda_i(A)+\lambda_j(B).
\]
\end{theorem}
\begin{proof}
Let $U$ be spanned by the eigenvectors of $A$ for $\lambda_i(A),\dots,\lambda_n(A)$, and $U'$ by those of $B$ for
$\lambda_j(B),\dots,\lambda_n(B)$. Then $\dim U=n-i+1$, $\dim U'=n-j+1$, and $\dim(U\cap U')\ge n-i-j+2=n-k+1$ with
$k=i+j-1$. On $U$, $R_A\le\lambda_i(A)$, and on $U'$, $R_B\le\lambda_j(B)$, by~\eqref{linear:eq:rayleigh}. So on a subspace $W$ of
$U\cap U'$ of dimension $n-k+1$, $R_{A+B}=R_A+R_B\le\lambda_i(A)+\lambda_j(B)$, and the min-max formula of
Theorem~\ref{linear:thm:courantfischer} gives $\lambda_k(A+B)\le\max_WR_{A+B}\le\lambda_i(A)+\lambda_j(B)$.
\end{proof}

\begin{corollary}[stability of eigenvalues]\label{linear:cor:weylpert}
For Hermitian $A,B$ and every $k$, $\lambda_n(B)\le\lambda_k(A+B)-\lambda_k(A)\le\lambda_1(B)$. In particular
$|\lambda_k(A+B)-\lambda_k(A)|\le\|B\|$, where $\|B\|=\max(|\lambda_1(B)|,|\lambda_n(B)|)$ is the operator norm of $B$.
\end{corollary}
\begin{proof}
Weyl with $j=1$ gives $\lambda_k(A+B)\le\lambda_k(A)+\lambda_1(B)$. Applied to $A+B$ and $-B$, it gives $\lambda_k(A)\le
\lambda_k(A+B)+\lambda_1(-B)=\lambda_k(A+B)-\lambda_n(B)$. For the operator norm, $\|Bx\|^2=\sum_i\lambda_i(B)^2|c_i|^2\le
\max_i\lambda_i(B)^2\|x\|^2$ in an eigenbasis, with equality on an eigenvector.
\end{proof}

Contrast with Chapter~\ref{linear:ch:eigen}, where a perturbation of size $10^{-6}$ moved the eigenvalues of a $12\times12$ Jordan
block by about $0.3$ (Figure~\ref{linear:fig:perturb}). Hermitian matrices are perfectly stable: eigenvalues move no more than the matrix does.
This is why numerical analysts love symmetric problems (Exercise~\ref{linear:ex:nonnormal} compares the two worlds).

\begin{remark}
Weyl's inequalities are only the first of many. In 1962 Alfred Horn conjectured a complete list of linear inequalities
describing all possible triples of spectra $(\lambda(A),\lambda(B),\lambda(A+B))$, defined recursively in terms of smaller
triples. The conjecture was proved in 1998--1999 by Klyachko, using geometric invariant theory, and by Knutson and Tao,
whose ``honeycombs'' proved the saturation property it needed. The same inequalities govern tensor products of
representations of $GL_n$ and intersections of Schubert varieties: see Fulton's survey in the sources.
\end{remark}

\section{Diagonals and spectra: Schur and Horn}\label{linear:sec:schurhorn}

For $x\in\R^n$, write $x^\downarrow$ for its entries in decreasing order. We say that $x$ is \emph{majorised} by $y$, written
$x\prec y$, if
\[
\sum_{i=1}^kx^\downarrow_i\le\sum_{i=1}^ky^\downarrow_i\quad(1\le k\le n),\qquad\text{with equality for }k=n.
\]
Informally, the entries of $y$ are more spread out than those of $x$: $y$ can be flattened into $x$ by moving mass from
large entries to small ones. For instance $(\frac13,\frac13,\frac13)\prec x\prec(1,0,0)$ for every probability vector $x$. A matrix with
non-negative entries whose rows and columns all sum to $1$ is \emph{doubly stochastic}; permutation matrices are the
simplest examples. We need one lemma about them.

\begin{lemma}\label{linear:lem:birkhoff}
Let $S$ be doubly stochastic and $(c_{ij})$ real numbers. There is a permutation $\sigma$ with $\sum_ic_{i\sigma(i)}\le
\sum_{i,j}c_{ij}s_{ij}$.
\end{lemma}
\begin{proof}
Call an entry of $S$ fractional if it lies strictly between $0$ and $1$. If there is none, $S$ is a permutation matrix and
we are done. Otherwise, each row containing a fractional entry contains at least two (its entries sum to $1$), and so does
each column. In the bipartite graph whose vertices are the rows and the columns, with an edge $\{i,j\}$ for each fractional
$s_{ij}$, every non-isolated vertex has degree at least $2$, so walking along edges without immediately going back must
eventually revisit a vertex: there is a cycle, which alternates between rows and columns and has even length. Add $\varepsilon$
and $-\varepsilon$ alternately to the entries of $S$ along the cycle. Every row and column of the cycle receives one $+\varepsilon$
and one $-\varepsilon$, so the sums are unchanged. The quantity $f(S)=\sum c_{ij}s_{ij}$ changes by $\varepsilon\delta$ for a
fixed $\delta$; choose the sign of $\varepsilon$ with $\varepsilon\delta\le0$, and increase $|\varepsilon|$ until some entry of
the cycle reaches $0$ or $1$. The new matrix is doubly stochastic, $f$ has not increased, and it has fewer fractional entries.
Repeating, we reach a permutation matrix.
\end{proof}

(Lemma~\ref{linear:lem:birkhoff} is the heart of Birkhoff's theorem: the doubly stochastic matrices are exactly the convex
combinations of permutation matrices. We will only need the lemma.)

\begin{theorem}[Schur, 1923]\label{linear:thm:schurmaj}
The diagonal of a Hermitian matrix is majorised by its eigenvalues: $(a_{11},\dots,a_{nn})\prec(\lambda_1,\dots,\lambda_n)$.
\end{theorem}
\begin{proof}
Write $A=UDU^*$. Then $a_{ii}=\sum_j|u_{ij}|^2\lambda_j$, and the matrix $(|u_{ij}|^2)$ is doubly stochastic because the rows
and columns of $U$ are unit vectors. The traces agree, which is the case $k=n$. Let $I$ be the set of indices of the $k$
largest diagonal entries, and $t_j=\sum_{i\in I}|u_{ij}|^2$. Then $0\le t_j\le1$, $\sum_jt_j=k$, and
\[
\sum_{j\le k}\lambda_j-\sum_{i\in I}a_{ii}=\sum_{j\le k}(1-t_j)\lambda_j-\sum_{j>k}t_j\lambda_j
\ge\lambda_k\sum_{j\le k}(1-t_j)-\lambda_k\sum_{j>k}t_j=\lambda_k\Bigl(k-\sum_jt_j\Bigr)=0,
\]
using $\lambda_j\ge\lambda_k$ for $j\le k$ and $\lambda_j\le\lambda_k$ for $j>k$.
\end{proof}

Applied in any orthonormal basis, Schur's theorem says that for orthonormal vectors $x_1,\dots,x_k$,
$\sum_{i\le k}\langle Ax_i,x_i\rangle\le\lambda_1+\dots+\lambda_k$, with equality for eigenvectors: \emph{Ky Fan's maximum
principle} (Exercise~\ref{linear:ex:kyfan}). The converse of Schur's theorem is the surprise of this section.

\begin{theorem}[Horn, 1954]\label{linear:thm:horn}
If $d\prec\lambda$ in $\R^n$, there is a real symmetric matrix with diagonal $d$ and eigenvalues $\lambda$.
\end{theorem}
\begin{proof}
We show that there is a real orthogonal $Q$ such that $Q\Lambda Q^{\mathsf T}$, with $\Lambda=\operatorname{diag}(\lambda)$,
has diagonal $d$. Conjugating by permutation matrices reorders diagonals and spectra, so we may assume that $d$ and
$\lambda$ are both decreasing. We argue by induction on $n$, the case $n=1$ being $d_1=\lambda_1$.

Let $n\ge2$. We have $\lambda_1\ge d_1$ (majorisation with $k=1$), and $d_1\ge\frac1n\sum d_i=\frac1n\sum\lambda_i\ge
\lambda_n$, so there is an index $k\le n-1$ with $\lambda_k\ge d_1\ge\lambda_{k+1}$. Choose $c\in[0,1]$ with
$c^2\lambda_k+(1-c^2)\lambda_{k+1}=d_1$, put $s=\sqrt{1-c^2}$, and conjugate the block $\operatorname{diag}(\lambda_k,
\lambda_{k+1})$ by the rotation $\bigl(\begin{smallmatrix}c&-s\\s&c\end{smallmatrix}\bigr)$:
\[
\begin{pmatrix}c&-s\\s&c\end{pmatrix}\begin{pmatrix}\lambda_k&0\\0&\lambda_{k+1}\end{pmatrix}\begin{pmatrix}c&s\\-s&c\end{pmatrix}
=\begin{pmatrix}d_1&cs(\lambda_k-\lambda_{k+1})\\cs(\lambda_k-\lambda_{k+1})&\lambda_k+\lambda_{k+1}-d_1\end{pmatrix}.
\]
After moving index $k$ to the front, $\Lambda$ is orthogonally similar to $M=\bigl(\begin{smallmatrix}d_1&r^{\mathsf T}\\r&
\Lambda'\end{smallmatrix}\bigr)$, where $\Lambda'$ is diagonal with entries
\[
\lambda'=(\lambda_1,\dots,\lambda_{k-1},\ \lambda_k+\lambda_{k+1}-d_1,\ \lambda_{k+2},\dots,\lambda_n).
\]
This vector is decreasing, because its $k$-th entry lies in $[\lambda_{k+1},\lambda_k]$. We claim $(d_2,\dots,d_n)\prec
\lambda'$. The totals agree. For $m\le k-1$, the sum of the first $m$ entries of $\lambda'$ is $\sum_{i\le m}\lambda_i\ge
m\lambda_k\ge md_1\ge\sum_{i=2}^{m+1}d_i$. For $m\ge k$, it is $\sum_{i\le m+1}\lambda_i-d_1\ge\sum_{i\le m+1}d_i-d_1=
\sum_{i=2}^{m+1}d_i$. By induction there is an orthogonal $Q'$ with $Q'\Lambda'Q'^{\mathsf T}$ having diagonal $(d_2,\dots,
d_n)$. Then $\operatorname{diag}(1,Q')\,M\,\operatorname{diag}(1,Q'^{\mathsf T})=\bigl(\begin{smallmatrix}d_1&(Q'r)^{\mathsf
T}\\Q'r&Q'\Lambda'Q'^{\mathsf T}\end{smallmatrix}\bigr)$ has diagonal $d$ and is orthogonally similar to $\Lambda$.
\end{proof}

So the second question of the opening has answer yes: $(1,1,1)\prec(3,0,0)$, and indeed the all-ones matrix $J$ has
diagonal $(1,1,1)$ and eigenvalues $3,0,0$. The proof is an algorithm, $n-1$ rotations, and a numerical calculation runs it.

\section{How far eigenvalues move}

Corollary~\ref{linear:cor:weylpert} bounds each eigenvalue separately by the operator norm. Measuring the perturbation with the
Frobenius norm $\|M\|_F=\bigl(\sum_{i,j}|m_{ij}|^2\bigr)^{1/2}=\sqrt{\operatorname{tr}M^*M}$, which is unitarily invariant, gives
a bound for all eigenvalues at once, and for normal matrices too.

\begin{theorem}[Hoffman--Wielandt, 1953]\label{linear:thm:hoffmanwielandt}
Let $A$ and $B$ be normal with eigenvalues $\alpha_1,\dots,\alpha_n$ and $\beta_1,\dots,\beta_n$. There is a permutation $\sigma$
with
\[
\sum_{i=1}^n|\alpha_i-\beta_{\sigma(i)}|^2\le\|A-B\|_F^2.
\]
If $A$ and $B$ are Hermitian, one can take the eigenvalues of both in decreasing order: $\sum_i(\lambda_i(A)-\lambda_i(B))^2
\le\|A-B\|_F^2$.
\end{theorem}
\begin{proof}
Write $A=UD_\alpha U^*$ and $B=VD_\beta V^*$, and $W=U^*V$, unitary. By unitary invariance,
\[
\|A-B\|_F^2=\|D_\alpha W-WD_\beta\|_F^2=\sum_{i,j}|\alpha_i-\beta_j|^2|w_{ij}|^2,
\]
since the $(i,j)$ entry of $D_\alpha W-WD_\beta$ is $(\alpha_i-\beta_j)w_{ij}$. The matrix $(|w_{ij}|^2)$ is doubly stochastic,
and Lemma~\ref{linear:lem:birkhoff} with $c_{ij}=|\alpha_i-\beta_j|^2$ gives the permutation. In the Hermitian case, if
$\alpha_i>\alpha_l$ but $\beta_{\sigma(i)}<\beta_{\sigma(l)}$, swapping $\sigma(i)$ and $\sigma(l)$ changes the sum by
\[
(\alpha_i-\beta_{\sigma(l)})^2+(\alpha_l-\beta_{\sigma(i)})^2-(\alpha_i-\beta_{\sigma(i)})^2-(\alpha_l-\beta_{\sigma(l)})^2
=-2(\alpha_i-\alpha_l)(\beta_{\sigma(l)}-\beta_{\sigma(i)})<0,
\]
so a minimising permutation can be taken order-preserving.
\end{proof}

\begin{beauty}
The eigenvalues of a Hermitian matrix are defined by a polynomial equation, yet everything in this chapter came from
optimisation: an eigenvalue is a max of mins of an average. From that one idea, interlacing, stability and convexity follow
in a few lines each, and Schur's and Horn's theorems say that the diagonals compatible with a spectrum form exactly the
convex polytope cut out by majorisation, the permutohedron (Rado's theorem identifies it with the convex hull of the
permutations of the spectrum). The analysis of eigenvalues is geometry in disguise.
\end{beauty}

\section{Exercises}

\begin{exercise}[$\star$]\label{linear:ex:normaliff}
Show that $A$ is normal if and only if $\|Ax\|=\|A^*x\|$ for all $x$. Deduce directly that an eigenvector of a normal $A$
for $\lambda$ is an eigenvector of $A^*$ for $\bar\lambda$.
\end{exercise}

\begin{exercise}[$\star$]\label{linear:ex:hornsmall}
Find a real symmetric matrix with diagonal $(2,1,0)$ and eigenvalues $(3,0,0)$, by following the proof of
Theorem~\ref{linear:thm:horn}. Is there one with diagonal $(2,2,-1)$ and eigenvalues $(3,0,0)$?
\end{exercise}

\begin{exercise}[$\star$]\label{linear:ex:pathgraph}
Let $A_n$ be the adjacency matrix of the path with $n$ vertices (ones just above and below the diagonal). Show that the
vectors $v_k=\bigl(\sin\frac{jk\pi}{n+1}\bigr)_{j=1}^n$ are eigenvectors, with eigenvalues $2\cos\frac{k\pi}{n+1}$, and check
the interlacing of Theorem~\ref{linear:thm:interlacing} between $A_{n-1}$ and $A_n$.
\end{exercise}

\begin{exercise}[$\star\star$]\label{linear:ex:degree}
Let $A$ be the adjacency matrix of a graph with average degree $\bar d$ and maximum degree $\Delta$. Show that $\bar d\le
\lambda_1(A)\le\Delta$, with equality on the left if and only if the graph is regular.
\end{exercise}

\begin{exercise}[$\star$]\label{linear:ex:nonnormal}
For $A=\bigl(\begin{smallmatrix}0&1\\0&0\end{smallmatrix}\bigr)$ and $E=\bigl(\begin{smallmatrix}0&0\\\varepsilon&0
\end{smallmatrix}\bigr)$, $\varepsilon>0$, compute the eigenvalues of $A+E$. Show that both Corollary~\ref{linear:cor:weylpert} and
Theorem~\ref{linear:thm:hoffmanwielandt} fail for this pair when $\varepsilon<1$.
\end{exercise}

\begin{exercise}[$\star\star$]\label{linear:ex:rankone}
Let $A$ be Hermitian and $v$ a vector. Show that the eigenvalues of $A+vv^*$ interlace those of $A$:
$\lambda_1(A+vv^*)\ge\lambda_1(A)\ge\lambda_2(A+vv^*)\ge\lambda_2(A)\ge\dots\ge\lambda_n(A+vv^*)\ge\lambda_n(A)$.
\end{exercise}

\begin{exercise}[$\star\star$]\label{linear:ex:kyfan}
(Ky Fan, 1949.) Show that $\lambda_1(A)+\dots+\lambda_k(A)=\max\sum_{i=1}^k\langle Ax_i,x_i\rangle$ over orthonormal families
$(x_1,\dots,x_k)$. Deduce that $\lambda_1+\dots+\lambda_k$ is a convex function of $A$, and that
$\sum_{i\le k}\lambda_i(A+B)\le\sum_{i\le k}\lambda_i(A)+\sum_{i\le k}\lambda_i(B)$.
\end{exercise}

\begin{exercise}[$\star\star$]\label{linear:ex:sqrtpsd}
Show that a positive semidefinite matrix $A$ has a unique positive semidefinite square root $\sqrt A$, and that $\sqrt A$
is a polynomial in $A$.
\end{exercise}

\begin{exercise}[$\star\star$]\label{linear:ex:simultaneous}
Show that commuting Hermitian matrices $A$ and $B$ have a common orthonormal eigenbasis. Is this true for commuting normal
matrices?
\end{exercise}

\begin{exercise}[$\star\star$]\label{linear:ex:realnormal}
Let $A$ be a real normal matrix. Show that there is a real orthogonal $Q$ such that $Q^{\mathsf T}AQ$ is block diagonal
with blocks $(a)$ of size $1$ and $\bigl(\begin{smallmatrix}a&-b\\b&a\end{smallmatrix}\bigr)$ of size $2$. What does this say
about a rotation of $\R^3$?
\end{exercise}

\begin{exercise}[$\star\star\star$]\label{linear:ex:bordered}
Let $\lambda_1>\mu_1>\lambda_2>\mu_2>\dots>\mu_{n-1}>\lambda_n$. Show that there are real $b_1,\dots,b_{n-1}$ and $c$ such that
the matrix $\bigl(\begin{smallmatrix}D&b\\b^{\mathsf T}&c\end{smallmatrix}\bigr)$, with $D=\operatorname{diag}(\mu_1,\dots,
\mu_{n-1})$, has eigenvalues $\lambda_1,\dots,\lambda_n$. (This is the converse of Cauchy's interlacing.)
\end{exercise}

\begin{exercise}[$\star\star\star$]\label{linear:ex:vonneumann}
(von Neumann's trace inequality, Hermitian case.) Show that $\operatorname{tr}(AB)\le\sum_i\lambda_i(A)\lambda_i(B)$ for
Hermitian $A,B$. When is there equality?
\end{exercise}

\section{Hints}

\begin{hint}{linear:ex:normaliff}
$\|Ax\|^2-\|A^*x\|^2=\langle(A^*A-AA^*)x,x\rangle$, and a Hermitian operator with all Rayleigh quotients zero is zero. Then
apply the first part to $A-\lambda I$.
\end{hint}
\begin{hint}{linear:ex:hornsmall}
Here $\lambda_1=3\ge d_1=2\ge\lambda_2=0$. For the second, check the partial sums.
\end{hint}
\begin{hint}{linear:ex:pathgraph}
$\sin(j-1)\theta+\sin(j+1)\theta=2\cos\theta\sin j\theta$, and $\sin0=\sin(n+1)\theta_k=0$.
\end{hint}
\begin{hint}{linear:ex:degree}
Use the Rayleigh quotient of the all-ones vector, and look at the largest coordinate of an eigenvector.
\end{hint}
\begin{hint}{linear:ex:nonnormal}
The characteristic polynomial is $X^2-\varepsilon$.
\end{hint}
\begin{hint}{linear:ex:rankone}
Weyl's inequalities with $B=vv^*$, whose eigenvalues are $\|v\|^2,0,\dots,0$.
\end{hint}
\begin{hint}{linear:ex:kyfan}
Extend the family to an orthonormal basis and apply Theorem~\ref{linear:thm:schurmaj} to the matrix of $A$ in that basis.
\end{hint}
\begin{hint}{linear:ex:sqrtpsd}
A square root commutes with $A$, so it preserves the eigenspaces of $A$.
\end{hint}
\begin{hint}{linear:ex:simultaneous}
$B$ preserves each eigenspace of $A$ and is Hermitian there.
\end{hint}
\begin{hint}{linear:ex:realnormal}
For a non-real eigenvalue $a+ib$ with eigenvector $x+iy$, show that $x\perp y$, $\|x\|=\|y\|$, and that $\operatorname{span}(x,y)$
and its orthogonal are invariant.
\end{hint}
\begin{hint}{linear:ex:bordered}
Expand $\det(XI-M)=(X-c)\prod_i(X-\mu_i)-\sum_ib_i^2\prod_{j\ne i}(X-\mu_j)$ and use partial fractions for
$\prod_j(X-\lambda_j)/\prod_i(X-\mu_i)$.
\end{hint}
\begin{hint}{linear:ex:vonneumann}
Diagonalise both, as in the proof of Theorem~\ref{linear:thm:hoffmanwielandt}, and use Lemma~\ref{linear:lem:birkhoff} and a swapping argument.
\end{hint}

\section{Solutions}

\begin{solution}{linear:ex:normaliff}
\idea A Hermitian operator is determined by its quadratic form, so ``same norms'' should mean ``same $A^*A$ and $AA^*$''.
\writeup $\|Ax\|^2-\|A^*x\|^2=\langle A^*Ax,x\rangle-\langle AA^*x,x\rangle=\langle Hx,x\rangle$ with $H=A^*A-AA^*$ Hermitian.
If $A$ is normal, this vanishes. Conversely, if $\langle Hx,x\rangle=0$ for all $x$, then all eigenvalues of $H$ are zero
(take $x$ an eigenvector), and $H=0$ by the spectral theorem. If $A$ is normal, so is $A-\lambda I$, whose adjoint is
$A^*-\bar\lambda I$; applying the first part to it, $Ax=\lambda x$ gives $\|(A^*-\bar\lambda)x\|=\|(A-\lambda)x\|=0$.
\end{solution}

\begin{solution}{linear:ex:hornsmall}
\idea Run the induction of Horn's proof by hand: one rotation places $d_1$, and the rest is a $2\times2$ problem.
\writeup With $\lambda=(3,0,0)$ and $d=(2,1,0)$: $k=1$, and $3c^2=2$ gives $c=\sqrt{2/3}$, $s=\sqrt{1/3}$. Conjugating
$\operatorname{diag}(3,0)$ by the rotation gives $\bigl(\begin{smallmatrix}2&\sqrt2\\\sqrt2&1\end{smallmatrix}\bigr)$, so
$\Lambda$ is orthogonally similar to
\[
M=\begin{pmatrix}2&\sqrt2&0\\\sqrt2&1&0\\0&0&0\end{pmatrix},
\]
and $\lambda'=(1,0)$ must be turned into $(d_2,d_3)=(1,0)$: nothing to do. So $M$ itself works: its diagonal is $(2,1,0)$
and, being $\operatorname{diag}(\text{rotated }2\times2,0)$, it has eigenvalues $3,0,0$. For diagonal $(2,2,-1)$: the
partial sums are $2,4,3$ against $3,3,3$, and $4>3$, so there is no such matrix, by Theorem~\ref{linear:thm:schurmaj}.
\end{solution}

\begin{solution}{linear:ex:pathgraph}
\idea Eigenvectors of a second-difference operator are sampled sines, the discrete vibrating string.
\writeup With $\theta=\frac{k\pi}{n+1}$, $(A_nv_k)_j=\sin(j-1)\theta+\sin(j+1)\theta=2\cos\theta\sin j\theta$ for $1\le j\le n$,
where the terms $\sin0\cdot\theta$ and $\sin(n+1)\theta$ that appear at the ends are $0$. The $n$ values $2\cos\frac{k\pi}
{n+1}$, $k=1,\dots,n$, are distinct, so they are all the eigenvalues. $A_{n-1}$ is the top-left corner of $A_n$, and the
interlacing $\cos\frac{k\pi}{n+1}\ge\cos\frac{k\pi}n\ge\cos\frac{(k+1)\pi}{n+1}$ holds because $\frac k{n+1}\le\frac kn\le
\frac{k+1}{n+1}$ and cosine decreases on $[0,\pi]$.
\end{solution}

\begin{solution}{linear:ex:degree}
\idea A lower bound for $\lambda_1$ needs one good test vector; an upper bound needs control of an eigenvector, best at its
largest coordinate.
\writeup With $\mathbf 1$ the all-ones vector, $R_A(\mathbf 1)=\frac{\sum_{i,j}a_{ij}}n=\frac{2|E|}n=\bar d$, so
$\lambda_1\ge\bar d$, with equality if and only if $\mathbf 1$ is an eigenvector for $\lambda_1$ (equality in a Rayleigh
bound happens only on the top eigenspace, by~\eqref{linear:eq:rayleigh}), that is, all degrees are equal. For the upper bound, let
$Ax=\lambda_1x$ and $|x_i|$ be the largest coordinate. Then $\lambda_1|x_i|=\bigl|\sum_{j\sim i}x_j\bigr|\le\deg(i)|x_i|\le
\Delta|x_i|$.
\end{solution}

\begin{solution}{linear:ex:nonnormal}
\idea The Jordan block of Chapter~\ref{linear:ch:eigen} again, now as a counterexample.
\writeup $\chi_{A+E}=X^2-\varepsilon$, so the eigenvalues are $\pm\sqrt\varepsilon$, while those of $A$ are $0,0$. The
perturbation has operator norm $\varepsilon$ and Frobenius norm $\varepsilon$, but the eigenvalues moved by $\sqrt\varepsilon>
\varepsilon$, and $\sum|\alpha_i-\beta_{\sigma(i)}|^2=2\varepsilon>\varepsilon^2$. Normality (here, of $A$) is essential.
\end{solution}

\begin{solution}{linear:ex:rankone}
\idea Adding a positive semidefinite matrix can only push eigenvalues up, and a rank-one one cannot push any of them past
the previous one.
\writeup $B=vv^*$ has eigenvalues $\|v\|^2,0,\dots,0$. Corollary~\ref{linear:cor:weylpert} gives $\lambda_k(A+B)\ge\lambda_k(A)+
\lambda_n(B)=\lambda_k(A)$. Weyl's inequality with $i=k-1$, $j=2$ gives $\lambda_k(A+B)\le\lambda_{k-1}(A)+\lambda_2(B)=
\lambda_{k-1}(A)$ for $k\ge2$.
\end{solution}

\begin{solution}{linear:ex:kyfan}
\idea The diagonal of $A$ in a basis starting with $x_1,\dots,x_k$ contains the numbers $\langle Ax_i,x_i\rangle$, and Schur
bounds the sum of any $k$ diagonal entries.
\writeup Extend to an orthonormal basis $x_1,\dots,x_n$; the matrix of $A$ in it is Hermitian with the same eigenvalues and
diagonal entries $\langle Ax_i,x_i\rangle$. The sum of any $k$ diagonal entries is at most the sum of the $k$ largest, which
is at most $\lambda_1+\dots+\lambda_k$ by Theorem~\ref{linear:thm:schurmaj}. Eigenvectors give equality. A maximum of linear
functions of $A$ is convex, and $\max(f+g)\le\max f+\max g$ gives the last inequality.
\end{solution}

\begin{solution}{linear:ex:sqrtpsd}
\idea Existence is the functional calculus; for uniqueness, anything commuting with $A$ lives on its eigenspaces.
\writeup Write $A=\sum_\lambda\lambda P_\lambda$ and set $\sqrt A=\sum_\lambda\sqrt\lambda P_\lambda$, positive semidefinite.
By Lagrange interpolation there is a polynomial $p$ with $p(\lambda)=\sqrt\lambda$ at the eigenvalues, and $p(A)=\sqrt A$.
If $S\succeq0$ and $S^2=A$, then $S$ commutes with $A=S^2$, so $S$ preserves each eigenspace $E_\lambda$ of $A$. On
$E_\lambda$, $S$ is Hermitian with eigenvalues $\mu\ge0$ satisfying $\mu^2=\lambda$, so $\mu=\sqrt\lambda$ and $S=\sqrt\lambda I$
there. Hence $S=\sqrt A$.
\end{solution}

\begin{solution}{linear:ex:simultaneous}
\idea Diagonalise $A$ first; $B$ cannot mix its eigenspaces.
\writeup If $Ax=\lambda x$, then $A(Bx)=BAx=\lambda Bx$, so $B$ preserves $E_\lambda(A)$, and $B|_{E_\lambda}$ is Hermitian.
Choose an orthonormal eigenbasis of $B|_{E_\lambda}$ in each eigenspace of $A$ and put them together: the eigenspaces are
mutually orthogonal, so this is an orthonormal basis of eigenvectors of both. For commuting normal matrices the same
argument works, since $B|_{E_\lambda}$ is normal (its adjoint is $B^*|_{E_\lambda}$, because $B^*$ also commutes with $A$:
$B^*$ is a polynomial in $B$, by the spectral decomposition and Lagrange interpolation, as in the previous solution).
\end{solution}

\begin{solution}{linear:ex:realnormal}
\idea Complexify, use the complex spectral theorem, and pair each non-real eigenvector with its conjugate.
\writeup Real eigenvalues give real eigenvectors (solve a real linear system). Let $\lambda=a+ib$ with $b\ne0$ and
$A(x+iy)=\lambda(x+iy)$, with $x,y$ real. Taking real and imaginary parts, $Ax=ax-by$ and $Ay=bx+ay$, so $P=\operatorname{span}
(x,y)$ is invariant and $A$ acts on it, in the basis $(x,y)$ up to scaling, by $\bigl(\begin{smallmatrix}a&b\\-b&a
\end{smallmatrix}\bigr)$. Since $A$ is real, $x-iy$ is an eigenvector for $\bar\lambda\ne\lambda$, so it is orthogonal to
$x+iy$ (Corollary~\ref{linear:cor:hermitian}): $\langle x+iy,x-iy\rangle=\|x\|^2-\|y\|^2+2i\langle x,y\rangle=0$, so $\|x\|=\|y\|$
and $x\perp y$. Normalising and swapping $x,y$ if needed gives an orthonormal basis of $P$ with matrix $\bigl(\begin{smallmatrix}
a&-b\\b&a\end{smallmatrix}\bigr)$. By Exercise~\ref{linear:ex:normaliff}, $x\pm iy$ are also eigenvectors of $A^{\mathsf T}=A^*$, so
$A^{\mathsf T}$ preserves $P$ and $A$ preserves $P^\perp$; $A|_{P^\perp}$ is again normal, and induction finishes. A rotation
$Q$ of $\R^3$ ($Q^{\mathsf T}Q=I$, $\det Q=1$) has eigenvalues of modulus $1$, so its blocks are $\pm1$ and rotation blocks
$\bigl(\begin{smallmatrix}\cos\theta&-\sin\theta\\\sin\theta&\cos\theta\end{smallmatrix}\bigr)$; with $\det=1$ in dimension $3$
there is a block $(1)$: every rotation of space has an axis (Euler, 1776).
\end{solution}

\begin{solution}{linear:ex:bordered}
\idea The characteristic polynomial of a bordered diagonal matrix divided by $\prod(X-\mu_i)$ is ``$X-c$ minus a sum of
simple fractions'', so we must show that $f/g$ has that shape with positive numerators. Interlacing is exactly what makes the
signs right.
\writeup Let $f=\prod_j(X-\lambda_j)$ and $g=\prod_i(X-\mu_i)$. Expanding $\det(XI-M)$ along the last row and column (or by
the Schur complement $X-c-b^{\mathsf T}(XI-D)^{-1}b$, times $g$),
\[
\det(XI-M)=g(X)\Bigl(X-c-\sum_i\frac{b_i^2}{X-\mu_i}\Bigr).
\]
The partial fraction decomposition of $f/g$, with $\deg f=\deg g+1$ and distinct simple poles, is
$\frac fg=X-c_0-\sum_i\frac{r_i}{X-\mu_i}$ with $c_0=\sum_j\lambda_j-\sum_i\mu_i$ and $r_i=-\frac{f(\mu_i)}{g'(\mu_i)}$.
Now $f(\mu_i)=\prod_j(\mu_i-\lambda_j)$ has exactly $i$ negative factors ($j\le i$), and $g'(\mu_i)=\prod_{l\ne i}(\mu_i-\mu_l)$
has exactly $i-1$ ($l<i$). So $f(\mu_i)$ and $g'(\mu_i)$ have opposite signs and $r_i>0$. Take $c=c_0$ and $b_i=\sqrt{r_i}$:
then $\det(XI-M)=f$.
\end{solution}

\begin{solution}{linear:ex:vonneumann}
\idea Same computation as Hoffman--Wielandt: the trace is a doubly stochastic average of products $\alpha_i\beta_j$, and
a linear function of a doubly stochastic matrix is extremal at a permutation; then rearrange.
\writeup With $A=UD_\alpha U^*$, $B=VD_\beta V^*$ and $W=U^*V$, $\operatorname{tr}(AB)=\operatorname{tr}(D_\alpha WD_\beta
W^*)=\sum_{i,j}\alpha_i\beta_j|w_{ij}|^2$. Lemma~\ref{linear:lem:birkhoff} applied to $c_{ij}=-\alpha_i\beta_j$ gives $\sigma$ with
$\operatorname{tr}(AB)\le\sum_i\alpha_i\beta_{\sigma(i)}$. If $\alpha_i>\alpha_l$ and $\beta_{\sigma(i)}<\beta_{\sigma(l)}$,
swapping increases the sum by $(\alpha_i-\alpha_l)(\beta_{\sigma(l)}-\beta_{\sigma(i)})>0$, so the sorted pairing is the
largest: $\operatorname{tr}(AB)\le\sum_i\lambda_i(A)\lambda_i(B)$. Equality holds when $A$ and $B$ have a common orthonormal
eigenbasis ordering both spectra the same way (then $W$ is a permutation matching sorted to sorted); conversely one can show
that equality forces such a basis.
\end{solution}

\section*{Sources and further reading}

\begin{itemize}
\item A.-L.~Cauchy, \emph{Sur l'\'equation \`a l'aide de laquelle on d\'etermine les in\'egalit\'es s\'eculaires des
mouvements des plan\`etes}, Exercices de math\'ematiques 4 (1829): reality of the eigenvalues and interlacing.
\item J.~W.~Strutt, Lord Rayleigh, \emph{The Theory of Sound}, Macmillan, 1877.
\item D.~Hilbert, \emph{Grundz\"uge einer allgemeinen Theorie der linearen Integralgleichungen} (Erste Mitteilung), G\"ottinger
Nachrichten (1904), 49--91; L.~A.~Steen, \emph{Highlights in the history of spectral theory}, American Mathematical Monthly
80 (1973), 359--381, for the story and the quotation.
\item E.~Fischer, \emph{\"Uber quadratische Formen mit reellen Koeffizienten}, Monatshefte f\"ur Mathematik und Physik 16
(1905), 234--249; R.~Courant, \emph{\"Uber die Eigenwerte bei den Differentialgleichungen der mathematischen Physik},
Mathematische Zeitschrift 7 (1920), 1--57.
\item H.~Weyl, \emph{Das asymptotische Verteilungsgesetz der Eigenwerte linearer partieller Differentialgleichungen},
Mathematische Annalen 71 (1912), 441--479.
\item I.~Schur, \emph{\"Uber eine Klasse von Mittelbildungen mit Anwendungen auf die Determinantentheorie}, Sitzungsberichte
der Berliner Mathematischen Gesellschaft 22 (1923), 9--20; A.~Horn, \emph{Doubly stochastic matrices and the diagonal of a
rotation matrix}, American Journal of Mathematics 76 (1954), 620--630; N.~N.~Chan and K.-H.~Li, \emph{Diagonal elements and
eigenvalues of a real symmetric matrix}, Journal of Mathematical Analysis and Applications 91 (1983), 562--566 (the
rotation proof given here).
\item A.~J.~Hoffman and H.~W.~Wielandt, \emph{The variation of the spectrum of a normal matrix}, Duke Mathematical Journal 20
(1953), 37--39; G.~Birkhoff, \emph{Tres observaciones sobre el \'algebra lineal}, Revista Universidad Nacional de Tucum\'an,
Serie A, 5 (1946), 147--151.
\item K.~Fan, \emph{On a theorem of Weyl concerning eigenvalues of linear transformations I}, Proceedings of the National
Academy of Sciences 35 (1949), 652--655.
\item W.~Fulton, \emph{Eigenvalues, invariant factors, highest weights, and Schubert calculus}, Bulletin of the AMS 37 (2000),
209--249; A.~Knutson and T.~Tao, \emph{Honeycombs and sums of Hermitian matrices}, Notices of the AMS 48 (2001), 175--186.
\item R.~Bhatia, \emph{Matrix Analysis}, Springer, 1997, Chapters II--III and VI; R.~A.~Horn and C.~R.~Johnson, \emph{Matrix
Analysis}, 2nd ed., Cambridge University Press, 2013, Chapter~4.
\end{itemize}
